% Evolution of MP
% theory: financial vulnerabilities and stress 
% MP for mitigating risks
% framework
% methodology
% findings
% rules

\begin{abstract}
    In recent years macroprudential policy has been evolving, with the aim to more closely manage the build-up of cyclical vulnerabilities in the financial system. These systemic risks often arise from the financial system’s inherent natural tendency to accumulate vulnerabilities over the cycle, compounded by its structural characteristics that amplify those risks through strategic complementarities. When these risks materialize during periods of financial stress, they can jeopardize medium-term economic growth. Hence, proactively preventing and mitigating them through effective macroprudential tools is critical.
    This paper attempts to quantify these risks and the effects of macroprudential policy on them through the use of an empirical growth-at-risk (GaR) framework. Using Bayesian quantile regression, we assess how systemic risks impact GDP growth and explore the potential costs and benefits of macroprudential policy on its growth distribution. Our findings confirm that the any policy measure introduces a trade-off between mitigation of systemic risks on the tail of the growth distribution on one hand and preserving median GDP growth on the other. 
    We contribute to the existing literature by three main ways 1. we adopt a more flexible approach to modeling the effects of macroprudential policy, which takes into account the current level of vulnerabilities, stress level in the system, and lending capacity of the banks, 2. we estimate the long-term sustainable level of systemic risks relative to the resilience of the financial system, 3. we offer improvements to existing macroprudential policy rule frameworks in the current literature to augment the decision-making process. This rule shouldn't be viewed as a binding constraint, but rather as a guideline that could serve as a basis for preemptive action.
    \end{abstract}

%%Graphical abstract
%\begin{graphicalabstract}
%\includegraphics{grabs}
%\end{graphicalabstract}

%%Research highlights
%\begin{highlights}
%\item Research highlight 1
%\item Research highlight 2
%\end{highlights}

\begin{keyword}
%% keywords here, in the form: keyword \sep keyword, up to a maximum of 6 keywords
growth-at-risk \sep systemic risk \sep quantile regressions \sep macroprudential policy \sep policy stance
%% PACS codes here, in the form: \PACS code \sep code
%% MSC codes here, in the form: \MSC code \sep code
%% or \MSC[2008] code \sep code (2000 is the default)
\end{keyword}